
\section{Task Setting}

\label{sec:1tasksetting}
\begin{frame}[label={page:background}]{Background}
\begin{figure}
    \centering
    \includegraphics[width=0.8\linewidth]{Fig/Fig_LFCarFleetScheme.pdf}
    \label{fig:carfleet1dim}
\end{figure}
\begin{block}{Question 1}
    For each physical car, how to obtain {\color{blue}self-data} despite the {\color{red}modeling error} and {\color{red}measurement noise}?
\end{block}
\begin{block}{Question 2}
    For each follower car or the leader car, how to obtain the {\color{red}Real-Time data} of the {\color{blue}object car}?
\end{block}
\end{frame}

\begin{frame}{Task Setting}
    \begin{block}
        {Task 1}
        Use the {\color{blue}High-Gain/Linear Parameter Variable(HG/LPV) Method based observer}\footfullcite{zemouche2018high} or {\color{blue}Neural Network(NN) based observer }\footfullcite{zhuang2025sampled}to estimate the state of each observer.
    \end{block}
    \begin{block}
        {Task 2}
        Utilize the {\color{blue}Distribute Scheme}\footfullcite{xu2023distributed}\footfullcite{wu2021design} and the local observer in Task 1 to obtain the real-time data of the target car.
    \end{block}
\end{frame}

\begin{frame}[label={page:internplan}]{Internship Plan}
    \begin{figure}
        \centering
        \includegraphics[width=0.85\linewidth]{Fig/TaskTable.pdf}
        \caption{Tasks during the Internship}
        \label{fig:InternshipTask}
    \end{figure}
    The main idea here: Use the Distributed Observer Scheme to test two observer design methods on the UGV platform by simulation and real experiment.
\end{frame}

\section{Ongoing Tasks}
\begin{frame}{Simple Multi-Car Fleet}
    follower.py: Initiating the QLab and QCar, use the follower algorithm to connect cars one by one.
    \vskip 4mm
     vehicle\_control.py: Uploaded by ACC for controlling the leader car and plotting the control graph, which works with the inside RT model.
    \vskip 4mm
    \begin{figure}
        \centering
        \includegraphics[width=0.85\linewidth]{Fig/Fig_SimpleAchievement.pdf}
        \caption{Current Approach}
        \label{fig:enter-label}
    \end{figure}
    {\color{red}This structure is hard to modify and does not offer the API for reading/writing data.}
\end{frame}

\begin{frame}{To the General Case}
    \begin{figure}
        \centering
        \includegraphics[width=0.85\linewidth]{Fig/Fig_GeneralScheme.pdf}
        \caption{Improved General Approach}
        \label{fig:enter-label}
    \end{figure}
\end{frame}

\begin{frame}[label={page:MainMethod}]{CarFleet Class}
\label{sec:mainmethod}
\begin{figure}
    \centering
    \includegraphics[width=0.85\linewidth]{Fig/Fig_CarFleetClass.pdf}
    \label{fig:CarFleetClass}
\end{figure}
Tasks:
\begin{itemize}
    \item 1. Initiation: Connect the QLab, create, and position the QCars, etc.
    \item 2. Establish the leader-follower scheme, and create the APIs for different controllers and observers.
    \item 3. Establish the API for reading the data of all cars in the fleet.
\end{itemize}
\end{frame}




\section{Future Plan}


\label{sec:futureplan}
\begin{frame}[label={page:futureplangraph}]{Future Plan}

\begin{itemize}
    \item Improve the built General Leader-Follower Scheme and test in all simulation environments.  Built the Distributed Observer Scheme corresponding to the Leader-Follower Scheme. \vskip 10mm
    \item Apply the two Different Observer Techniques in the distributed observer scheme to test the performance of the UGV System. \vskip 10mm
    \item Test the build scheme in both Simulation and Realistic Environment.
\end{itemize}
\end{frame}

\begin{frame}{Ending}
    \begin{center}
        That's all. Thank you for listening!!
    \end{center}   
\end{frame}

